% ======================================================================
% Zero-Parameter Topological Mass Matrix Reconstruction and
% CKM/PMNS Closure from 88-Site Grid Dynamics
% IDCM v10.0 (It-From-Beat) — paper draft skeleton (2026-09-08)
% Status: DRAFT v0.1 — content mirrors 00_Master_Theorem.md
% ======================================================================
\documentclass[12pt]{article}
\usepackage[margin=1in]{geometry}
\usepackage{amsmath,amssymb,amsthm}
\usepackage{booktabs}
\usepackage{bm}
\usepackage[hidelinks]{hyperref}

\newtheorem{axiom}{Axiom}
\newtheorem{theorem}{Theorem}
\newtheorem{conjecture}{Conjecture}
\newtheorem{corollary}{Corollary}

\newcommand{\T}{\hat{T}}
\newcommand{\vev}[1]{\langle #1\rangle}

\title{\bf Zero-Parameter Topological Mass Matrix Reconstruction\\[2mm]
and CKM/PMNS Closure from 88-Site Grid Dynamics}
\author{IDCM v10.0 collaboration \\ \small (It-From-Beat program)}
\date{2026-09-08 \quad DRAFT v0.1}

\begin{document}
\maketitle

\begin{abstract}
We present a zero-free-parameter reconstruction of the Standard Model quark
and lepton mixing matrices from the discrete spectrum of an 88-site dual
chain (the $176$-Laplacian first 88 modes, $\gamma_k = 4\sin^2(k\pi/176)$).
All CKM matrix elements are reconstructed within $1\%$ (9/9 elements) from
four structural angles, with the suppressed sector $\{\sin\theta_{13},
|V_{ub}|\}$ generated by a two-rung ``k-ladder'' whose Fibonacci index
structure is traced to the transfer-matrix matrix elements
$F_n = (\T^n)_{12}$ of the tree sector ($\T = [[1,1],[1,0]]$), and whose
termination at exactly two rungs is forced by parity and budget existence
conditions. The PMNS matrix closes in the tree sector
$\mathbb{Q}(\sqrt{5})$ with $\delta_{CP} = 230^\circ$ exact and
$J_{CP} = 0.0255$ ($0.49\%$). We resolve the long-standing $|V_{ub}|$
tension as an outdated experimental anchor artefact. The Higgs VEV ratio
$\tan\beta \simeq L_4\cdot 43/L_5$ admits a trace-form mechanism
$\tan\beta = [\mathrm{Tr}(\Delta_{88})+\mathrm{Tr}(\T^{10})]/\mathrm{Tr}(\T^5)$
stated as an open theorem.
\end{abstract}

\section{Introduction}
\label{sec:intro}
The IDCM (Information Dynamics Cosmology Model) programme derives Standard
Model parameters from the spectrum of a discrete $88$-site dual chain
without free parameters \cite{o149,o155,o156}. This paper consolidates the
closure chain: (i) the beat-frequency identity
\eqref{eq:beatid}, (ii) the $k$-ladder theorem \eqref{eq:ladder} with
transfer-matrix provenance, (iii) the $|V_{ub}|$ closure
\eqref{eq:vub}, (iv) the 9/9 CKM reconstruction, and (v) the PMNS tree-sector
closure.

\section{Setup and Axioms}
\label{sec:setup}
\begin{axiom}[88-site dual chain]
The world is the first 88 modes of the 176-Laplacian:
$\gamma_k = 4\sin^2(k\pi/176)$, $k=1,\dots,88$. Even $k$ (world sector)
sum to $90$; odd $k$ (observer sector) sum to $88$; the total is
$178 = 2F_{11}$; the even--odd difference is
$\mu^2 = 2 = \gamma_{44}$.
\end{axiom}

\begin{axiom}[Beat operator]
$\omega_n = 2\sin(n\pi/88)$, $\mathrm{beat}(i,j) = |\omega_i - \omega_j|$.
Mixing quantities are beat ratios.
\end{axiom}

\begin{axiom}[Transfer matrix]
$\T = \begin{pmatrix}1&1\\1&0\end{pmatrix}$; Lucas $L_n = \mathrm{Tr}(\T^n)$;
Fibonacci $F_n = (\T^n)_{12}$.
\end{axiom}

\section{Beat-frequency theorems}
\label{sec:beat}
\begin{theorem}[Beat identity; O149]
\label{thm:beatid}
For all $a<b$,
\begin{equation}
\label{eq:beatid}
\mathrm{beat}(a,b) \;=\; 4\,\sin\!\Big(\frac{(b-a)\pi}{176}\Big)\,
v_g\!\Big(\frac{a+b}{2}\Big),
\qquad v_g(n) = \cos\frac{n\pi}{88}.
\end{equation}
\emph{Verified} $946/946$; equivalently $\gamma_n = 4v_g(44-n)^2$.
\end{theorem}

\begin{theorem}[$v_g$ floor; O149]
Adjacent-pair ratios satisfy
$\mathrm{beat}(n,n{+}1)/\mathrm{beat}(\cdot,\cdot) \ge
\sin(\pi/176)/\cos(3\pi/176) = 0.017875$. Since
$\sin\theta_{13} \approx 0.0035$ and $|V_{ub}| \approx 0.0037$ lie below this
floor, the two suppressed CKM elements necessarily require a two-factor
structure.
\end{theorem}

\section{The k-ladder theorem}
\label{sec:ladder}
\begin{theorem}[k-ladder; O149 + transfer-matrix form, this work]
\label{thm:ladder}
With $s_k = 44-(k{+}1)$ and $d_k = F_{k+6}$,
\begin{equation}
\label{eq:ladder}
g_k \;=\; \frac{\mathrm{beat}(44{-}k,\,45{-}k)}
{\mathrm{beat}\big(\tfrac{s_k+d_k}{2},\,\tfrac{s_k-d_k}{2}\big)}
\;=\; \frac{\sin\!\big(\tfrac{F_{k+2}\pi}{176}\big)\,
\sin\!\big(\tfrac{\pi}{176}\big)}
{\sin\!\big(\tfrac{F_{k+6}\pi}{176}\big)\,
\sin\!\big(\tfrac{(45+k)\pi}{176}\big)},
\end{equation}
where the closed form is exact ($1$e-25). The Fibonacci indices are
transfer-matrix matrix elements $F_n = (\T^n)_{12}$; the index gap is
$(k{+}6)-(k{+}2) = 4 = B$, the tree branching factor.
\end{theorem}

\begin{theorem}[Existence forcing]
\label{thm:existence}
The mode pair $((s_k\pm d_k)/2)$ must be positive integers $\le 44$.
$k=1$ fails parity ($s_1=42$ even vs $F_7=13$ odd $\Rightarrow$ half-integer
modes); $k=4$ fails the budget ($F_{10}=55 > s_4=39$). Hence
$k\in\{2,3\}$ is \emph{forced}: exactly the two suppressed CKM elements
$\sin\theta_{13}=g_2=0.0035074294$ and $|V_{ub}|=g_3=0.0036910176$.
\end{theorem}

\section{CKM closure}
\label{sec:ckm}
\begin{theorem}[$|V_{ub}|$ closure; this work]
\label{thm:vub}
\begin{equation}
\label{eq:vub}
|V_{ub}| = g_3 = 0.0036910176, \qquad
\theta_{13}^{\rm CKM} = \arcsin g_3 = 0.211480^\circ,
\end{equation}
agreeing with the PDG exclusive average $(3.70\pm0.11)\times10^{-3}$ at
$0.08\sigma$. The previously reported ``$1.7\sigma$ tension'' is traced to
an outdated anchor: $g_2 = 0.0035074294$ reproduces the
UTfit-2007-era average $|V_{ub}|=(3.51\pm0.19)\times10^{-3}$ at $0.01\sigma$,
not the modern value.
\end{theorem}

\begin{theorem}[9/9 CKM reconstruction]
\label{thm:ckm9}
The standard-parametrization matrix built from the four structural angles
$\{\theta_{12}{=}13.04^\circ,\ \theta_{23}{=}2.38^\circ,\
\theta_{13}{=}0.211480^\circ,\ \delta{=}66.399960^\circ\}$ reproduces all
nine CKM magnitudes within $1\%$ (max deviation $0.74\%$, $|V_{ub}|$
deviation $0.03\%$). The CP phase closes at
\begin{equation}
\label{eq:delta}
\delta_{CP} = \arctan\frac{\sin(21\pi/176)}{\sin(9\pi/176)}
= 66.39996005^\circ \qquad (1\text{e-}9),
\end{equation}
with $21+9 = 30 = N_d - 2L_4$ and $21-9 = 12 = B\,N_{\rm gen}$; the winding
$d=12 \equiv 0 \pmod 4$ satisfies the O155 parity condition.
\end{theorem}

\section{PMNS closure}
\label{sec:pmns}
\begin{theorem}[PMNS tree-sector closure]
\label{thm:pmns}
The PMNS matrix closes in the tree sector $\mathbb{Q}(\sqrt5)$:
\begin{align}
\theta_{12} &= \arctan\frac{33+\varphi}{33\varphi-1} && \text{exact}
\; (= \arctan\beta + 1/33),\\
\theta_{23} &= \arctan\frac{\sin(7\pi/55)}{\sin(6\pi/55)}
= 49.196588^\circ && (0.003\%, \text{ home } 55=F_{10},\ 33^2{+}44^2{=}55^2),\\
\theta_{13} &= \arcsin\!\big(\varepsilon\cdot\tfrac{32}{33}\big)
= 8.616868^\circ && (0.56\%),\\
\delta_{CP} &= 5\times46 = 230.00^\circ && \text{exact},\\
J_{CP} &= 0.025525 && (0.49\%).
\end{align}
The full matrix obeys $\max\Delta|U| = 0.0008$, $\mathrm{mean}\,\Delta|U|
= 0.0002$. The $1/33$ co-closure ($+1/33$ additive in $\theta_{12}$,
$\times32/33$ multiplicative in $\theta_{13}$) is the trace-density
decomposition $\mathrm{Tr}(P_B)/K_{\rm tree} = 1/33$,
$\mathrm{Tr}(P_I)/K_{\rm tree} = 32/33$, $P_B + P_I = \mathbf{1}$.
\end{theorem}

\section{The Higgs VEV ratio (open theorem)}
\label{sec:tanbeta}
\begin{conjecture}[tan$\beta$ trace theorem]
\label{cnj:tanb}
The semi-empirical VEV ratio $\tan\beta \simeq 27.3625$ admits the trace
form
\begin{equation}
\label{eq:tanb}
\tan\beta \;=\; \frac{\mathrm{Tr}(\Delta_{88})+\mathrm{Tr}(\T^{10})}
{\mathrm{Tr}(\T^{5})}
\;=\; \frac{178+123}{11} \;=\; \frac{L_4\cdot 43}{L_5}
\;=\; 27.363636\ldots \qquad (0.004\%),
\end{equation}
where $\mathrm{Tr}(\Delta_{88}) = 178 = 2F_{11}$ is the chain Laplacian
trace, $43 = B\,L_5-1$ is the boundary mode bracketing the $\mu^2$ site
$44$ (with $\omega_{43}=\omega_{45}$ degenerate), and the boundary mode is
anti-periodic ($n\cdot 44 \equiv 44 \pmod{88} \Rightarrow n$ odd). The
down-type VEV is consistent: $v_d = v\cos\beta = 8.992$ GeV
$\approx L_6/2 = 9 = N_{\rm gen}^2$ ($0.088\%$), while
$v_u = v\sin\beta \approx 2L_{10} = 246$ GeV.
\end{conjecture}
\textbf{Open statement:} prove that the geometric eigenstate ratio
$v_u/v_d$ necessarily projects onto the trace ratio
\eqref{eq:tanb}. Until then, $P7$ remains an open problem (candidate
unique at $0.005\%$ tolerance, mechanism package consistent but not a
theorem).

\section{Discussion}
\label{sec:discussion}
\begin{itemize}
\item \textbf{Number fields.} CKM lives in $\mathbb{Q}(\zeta_{176})^+$;
PMNS in $\mathbb{Q}(\sqrt5)$. The R12 domain theorem
($5\nmid 88 \Rightarrow \mathbb{Q}(\sqrt5)\cap\mathbb{Q}(\zeta_{88})^+
= \mathbb{Q}$) forces exact closures to be rational; no
$\mathbb{Q}(\zeta_{880})^+$ extension is admitted.
\item \textbf{Old-anchor artefact.} The $|V_{ub}|$ ``tension'' was a
comparison of a 2007-era-anchored value against the modern PDG average,
not a structural discrepancy; the k-ladder $k=3$ rung matches the modern
value at $0.08\sigma$.
\item \textbf{Null controls.} All locked quantities pass local-density
nulls (e.g.\ $\tan\beta$: 12 distinct hits at $0.1\%$, unique at
$0.005\%$; the k-ladder: 2 candidates out of 98 under relaxation,
$p_{\rm joint} = 9.6\times10^{-10}$).
\end{itemize}

\section{Open problems}
\label{sec:open}
\begin{enumerate}
\item \textbf{P7 ($\tan\beta$)}: prove Conjecture~\ref{cnj:tanb} (the
VEV-eigenstate projection onto the trace ratio).
\item \textbf{k-ladder index assignment}: derive why the numerator index is
$F_{k+2}$ and the denominator difference $F_{k+6}$ (gap $= B = 4$), and the
origin of the reflected index $45+k = 88-s_k$.
\item \textbf{Allocation rule}: elevate the channel-duality penalty
argument ($+1/33$ vs $\times 32/33$) to a theorem.
\item \textbf{Medium-sector angles}: 9/11 mixing angles remain
posterior; a k-ladder analogue in the unsuppressed sector is open.
\end{enumerate}

\begin{thebibliography}{9}
\bibitem{o149} O149, \emph{Group-velocity locking of suppressed CKM
elements}, IDCM v10.0 open-problems series, 2026-09-05.
\bibitem{o155} O155, \emph{K-tree eigen-operatorisation}, IDCM v10.0, 2026-09-08.
\bibitem{o156} O156, \emph{$\theta_{23}$ home scale 55 from first
principles}, IDCM v10.0, 2026-09-08.
\bibitem{o142} O142, \emph{Integer-constant tracking via
$\hat T$-traces}, IDCM v10.0, 2026-09-02.
\bibitem{o146} O146, \emph{Beat-existence audit and null-control},
IDCM v10.0, 2026-09-04.
\end{thebibliography}

\end{document}
